\chapter{Devices, Bias, and Regulation}
\label{ch:devices}
Passive circuits shape a signal; active devices supply the controlled current that
amplifies, switches, oscillates, and regulates it.  The useful mathematical model
is local: choose a DC operating point, linearize the device there, and use the
resulting controlled source with the feedback ideas of \cref{ch:feedback}.

\section{Bias Turns a Nonlinear Device into a Small-Signal Model}
\label{sec:devicebias}
A diode's current rises approximately exponentially with voltage,
\[
I_D=I_S\left(\ee^{V_D/(nV_T)}-1\right),
\]
where \(V_T\) is the thermal voltage (about \SI{26}{\milli\volt} near room
temperature) and \(n\) is an ideality factor. At a forward-bias point for which the
minus-one term is negligible, differentiating gives the small-signal resistance
\[
\boxed{\quad r_d\approx\frac{nV_T}{I_D}.\quad}
\]
This is a local incremental model. Large excursions, series resistance, junction
capacitance, temperature, and reverse behavior require a fuller device model.

A bipolar junction transistor (BJT) uses a forward-biased base--emitter junction
to control collector current.  Near its operating point,
\[
I_C\approx I_S\ee^{V_{BE}/V_T},\qquad
\boxed{\quad g_m=\frac{\partial I_C}{\partial V_{BE}}
\approx\frac{I_C}{V_T}.\quad}
\]
Thus, near the chosen bias point, a small \(v_{be}\) produces
\(i_c=g_mv_{be}\), and a collector load converts current to voltage. The base
current is approximately \(I_C/\beta\) in this local model. A field-effect
transistor similarly has \(i_d=g_mv_{gs}\) at its bias point. Its conductive gate
current is commonly much smaller than BJT base current, although leakage, protection
structures, and RF capacitive current still matter
\cite[pp.~123--130]{bowick2008rfcircuit}. Vacuum tubes can also be modeled as
voltage-controlled current sources, but the technologies differ in their detailed
device models.

\section{Load Lines and Emitter or Source Degeneration}
\label{sec:loadlines}
Consider a common-emitter BJT stage with collector resistor \(R_C\) and supply
\(V_{CC}\).  KVL constrains the possible DC points to
\[
V_{CE}=V_{CC}-I_CR_C.
\]
This straight \textbf{DC load line} intersects the transistor's nonlinear static
curves at the quiescent point. For the simplified resistively loaded Class-A stage,
bias near the middle of the available voltage swing permits approximately symmetric
excursion before cutoff or saturation. At RF, reactive matching and dynamic device
behavior make the static load-line picture mainly qualitative
\cite[pp.~458--460]{nguyen2015rfic}. Moving the bias toward cutoff underlies the
linearity/efficiency trade of Class B or C operation (\cref{sec:conductionangle}).

An emitter resistor \(R_E\) adds local negative feedback. If collector current rises,
emitter voltage rises, reducing \(V_{BE}\) and opposing the change. With large
\(\beta\), negligible transistor output resistance, an unloaded collector resistance,
and approximately equal emitter and collector incremental currents, the local model gives
\[
 r_e\approx\frac{1}{g_m}=\frac{V_T}{I_C},
 \qquad
\boxed{\quad A_v\approx-\frac{R_C}{r_e+R_E}.\quad}
\]
The minus sign is the common-emitter inversion. The resistor lowers gain but reduces
sensitivity to device parameters and raises input resistance; source degeneration gives
the analogous FET behavior \cite[pp.~665--667]{nguyen2015rfic}. Loading and finite
output resistance require the corresponding effective collector load in place of
\(R_C\).

\begin{workedbox}[: how a resistor stabilizes a BJT]
At \(I_C=\SI{1}{\milli\ampere}\), \(r_e\approx
\SI{26}{\milli\volt}/\SI{1}{\milli\ampere}=\SI{26}{\ohm}\).  With
\(R_C=\SI{4.7}{\kilo\ohm}\) and no unbypassed emitter resistor, the ideal unloaded local
gain is about \(-4700/26=-181\).  Adding \(R_E=\SI{470}{\ohm}\) changes it to
about \(-4700/(26+470)=-9.5\).  The stage gives up gain, but a modest change in
transconductance no longer makes its behavior unpredictable.  A bypass capacitor
can restore AC gain over a chosen band while retaining DC bias stabilization.
\end{workedbox}

\section{Rectification and Regulation Are Feedback Problems}
\label{sec:regulation}
A diode rectifier is deliberately used far outside its small-signal region: it
passes one polarity of an AC waveform to charge a reservoir capacitor. A full-wave
bridge replenishes the capacitor twice per line cycle, so
\(f_{\rm ripple}=2f_{\rm line}\) (\cref{ch:active}). If ripple is small, load current
is approximately constant between short recharge intervals, and transformer/source
impedance is neglected, charge balance gives
\[
\boxed{\quad\Delta V\approx\frac{I_{\rm load}}{f_{\rm ripple}C}.\quad}
\]
More capacitance, less load current, or more frequent charging reduces this approximate
ripple. A bridge has two forward-biased diodes in the conducting path, but their
forward drops depend on current and temperature rather than being fixed constants.

A voltage regulator senses output voltage, compares it with a reference, and adjusts
a pass element to reduce error. At frequencies where loop gain is large, feedback
reduces line- and load-induced output variation by approximately \(1+L\), as in
\cref{ch:feedback}. Stability depends on regulator topology and on poles and zeros
introduced by the pass device, output network, load, and compensation. Some regulators
are internally compensated; others require specified output capacitance, ESR, or load.
In every case the applicable data-sheet stability conditions take precedence over the
generic model.

\begin{exambox}
Bias, clipping, rectification, and regulation should not be collapsed into one
``diode/transistor'' fact.  Bias chooses a locally linear operating point;
rectification intentionally uses nonlinearity; regulation closes a feedback loop
around a source whose raw voltage would otherwise move.
\end{exambox}

\begin{mistakebox}
High input resistance at a FET gate does not mean every FET amplifier has high
gain, and a large BJT current gain \(\beta\) does not make its collector current
stable.  Gain and bias stability are circuit properties, set by the device \emph{and}
its surrounding impedances and feedback.
\end{mistakebox}
